\begin{lemma}[Lemma~3.1]
Let $n,m\in\mathbb N$, let $\sigma:\mathbb R^{n}\to\mathbb R^{n\times m}$,
$\sigma=[\sigma^1,\ldots,\sigma^m]$, be a frame of $m$ vector fields on $\mathbb R^{n}$,
and let $A:\mathbb R^{m}\to\mathbb R^{m\times m}$, $H:\mathbb R^{m}\to\mathbb R$ satisfy
the normalization (2.3) (\noderef{NormalizedAH}). Let $\psi\in C^{2}(\mathbb R)$, let
$w:\mathbb R^{n}\to\mathbb R$, let $x_0\in\mathbb R^{n}$ and let $(D,B)$ be a second-order
jet of $w$ at $x_0$ (\noderef{TwiceDiffAt}). Put
\[
 \xi_0:=\mathrm{hgrad}\ \sigma\ x_0\ D\quad(\noderef{hgrad}),\qquad
 t:=\psi'\big(w(x_0)\big),
\]
and assume $\xi_0\ne0$ and $t>0$. Let $D\otimes D$ denote the continuous bilinear form
$(h,k)\mapsto D(h)\,D(k)$. Then:
\begin{enumerate}
\item $\big(tD,\ tB+\psi''(w(x_0))\,D\otimes D\big)$ is a second-order jet of $\psi\circ w$
  at $x_0$;
\item $\mathrm{hgrad}\ \sigma\ x_0\ (tD)=t\,\xi_0$;
\item with that jet, (3.3) holds (\noderef{Lop}, \noderef{Ellip}):
\[
 \mathrm{Lop}\ \sigma\ A\ H\ x_0\ (tD)\ \big(tB+\psi''(w(x_0))\,D\otimes D\big)
 \ =\ t\Big[\mathrm{Lop}\ \sigma\ A\ H\ x_0\ D\ B
 -\frac{\psi''(w(x_0))}{\psi'(w(x_0))}\,\mathcal E\big(A,\xi_0\big)\Big].
\]
\end{enumerate}
\end{lemma}
\begin{proof}
Write $a:=w(x_0)$, $t=\psi'(a)$, $c:=\psi''(a)$, and $h:=y-x_0$ for the increment. Items
(1) and (2) are the paper's two displayed identities
$\mathcal X(\psi\circ w)=\psi'(w)\mathcal Xw=t\xi_0$ and
$\mathcal X\mathcal X^{s}(\psi\circ w)=\psi'(w)\mathcal X\mathcal X^{s}w+\psi''(w)\mathcal Xw\otimes\mathcal Xw$
at $x_0$, read on jets; item (3) is the computation of Lemma~3.1's proof, whose last step is
(2.5). Since (2.5) is a two-line consequence of (2.3) (pm-0012 keeps it
inline rather than as a node), we derive it in Step 4 below.

\medskip\noindent\textbf{Step 1: a Peano expansion for $\psi$ at $a$.} We claim: for every
$\varepsilon>0$ there is $\rho>0$ such that
\[
 \Big|\psi(a+s)-\psi(a)-t\,s-\tfrac12c\,s^{2}\Big|\ \le\ \varepsilon\,s^{2}
 \qquad\text{whenever }|s|<\rho. \tag{1}
\]
Indeed, $\psi''$ is continuous because $\psi\in C^{2}$, so there is $\rho>0$ with
$|\psi''(\theta)-\psi''(a)|\le2\varepsilon$ whenever $|\theta-a|\le\rho$. Let $|s|<\rho$.
If $s=0$ the claim is trivial. Otherwise, by Taylor's theorem with Lagrange remainder
applied to $\psi$ on the segment between $a$ and $a+s$, there is $\theta$ strictly between
$a$ and $a+s$ with
\[
 \psi(a+s)=\psi(a)+\psi'(a)s+\tfrac12\psi''(\theta)s^{2},
\]
and $|\theta-a|<|s|<\rho$, so
$|\psi(a+s)-\psi(a)-ts-\frac12cs^{2}|=\frac12|\psi''(\theta)-\psi''(a)|\,s^{2}\le\varepsilon s^{2}$.
This is (1). (Only $\psi\in C^{2}$ near $a$ is used; the hypothesis $t>0$ plays no role
here.)

\medskip\noindent\textbf{Step 2: proof of (1).} Set
\[
 D_\psi:=tD,\qquad B_\psi:=tB+c\,D\otimes D .
\]
$B_\psi$ is symmetric: $B$ is symmetric by \noderef{TwiceDiffAt} and
$(D\otimes D)(h,k)=D(h)D(k)=(D\otimes D)(k,h)$. It remains to prove the Peano estimate.

Write, for $y\in\mathbb R^{n}$ and $h=y-x_0$,
\[
 S(h):=w(y)-w(x_0),\qquad R(h):=S(h)-D(h)-\tfrac12B(h,h),
\]
so that by \noderef{TwiceDiffAt} applied to $w$: for every $\varepsilon_1>0$ there is
$\delta_1>0$ with
\[
 |R(h)|\le\varepsilon_1|h|^{2}\qquad\text{for }|h|<\delta_1. \tag{2}
\]
Fix once and for all $\delta_1^{0}>0$ obtained from (2) with $\varepsilon_1=1$, and put
\[
 C_0:=\|D\|+\tfrac12\|B\|+1 .
\]
For $|h|<\min\{\delta_1^{0},1\}$ we get from (2) with $\varepsilon_1=1$ that
\[
 |S(h)|\le|D(h)|+\tfrac12|B(h,h)|+|h|^{2}\le\big(\|D\|+\tfrac12\|B\|\,|h|+|h|\big)|h|
 \le C_0|h|, \tag{3}
\]
using $|h|\le1$. Likewise, for such $h$,
\[
 |S(h)-D(h)|\le\tfrac12\|B\|\,|h|^{2}+|h|^{2}=C_1|h|^{2},\qquad
 C_1:=\tfrac12\|B\|+1 , \tag{4}
\]
and $|S(h)+D(h)|\le(C_0+\|D\|)|h|$ by (3).

Now let $\varepsilon>0$ be given. Choose, in this order:
\begin{itemize}
\item $\varepsilon_2:=\dfrac{\varepsilon}{3(C_0^{2}+1)}>0$ and, by Step 1 with
  $\varepsilon_2$, a radius $\rho>0$ for which (1) holds;
\item $\varepsilon_1:=\dfrac{\varepsilon}{3(|t|+1)}>0$ and, by (2), a radius $\delta_1>0$
  for which (2) holds with this $\varepsilon_1$;
\item $\delta:=\min\Big\{\delta_1^{0},\,1,\,\delta_1,\,\dfrac{\rho}{C_0+1},\,
  \dfrac{2\varepsilon}{3\big(|c|+1\big)C_1\big(C_0+\|D\|+1\big)}\Big\}>0 .$
\end{itemize}
Let $|h|<\delta$. By (3) and $|h|<\rho/(C_0+1)$ we have $|S(h)|\le C_0|h|<\rho$, so (1)
applies with $s:=S(h)$. Decompose, using $\psi(w(y))=\psi(a+S(h))$ and
$\psi(w(x_0))=\psi(a)$:
\[
 \psi(w(y))-\psi(w(x_0))-D_\psi(h)-\tfrac12B_\psi(h,h)
 =\underbrace{\Big[\psi(a+S)-\psi(a)-tS-\tfrac12cS^{2}\Big]}_{=:T_1}
 +\underbrace{t\Big[S-D(h)-\tfrac12B(h,h)\Big]}_{=:T_2}
 +\underbrace{\tfrac12c\Big[S^{2}-D(h)^{2}\Big]}_{=:T_3},
\]
where $S=S(h)$; this is an identity, obtained by expanding
$D_\psi(h)=tD(h)$ and $\frac12B_\psi(h,h)=\frac t2B(h,h)+\frac c2D(h)^{2}$ and cancelling.
We estimate the three terms.

$|T_1|\le\varepsilon_2S^{2}\le\varepsilon_2C_0^{2}|h|^{2}\le\frac\varepsilon3|h|^{2}$ by
(1), (3) and the choice of $\varepsilon_2$.

$|T_2|=|t|\,|R(h)|\le|t|\varepsilon_1|h|^{2}\le\frac\varepsilon3|h|^{2}$ by (2) and the
choice of $\varepsilon_1$.

$|T_3|=\frac12|c|\,|S-D(h)|\,|S+D(h)|\le\frac12|c|\,C_1|h|^{2}(C_0+\|D\|)|h|
\le\frac\varepsilon3|h|^{2}$ by (4), (3) and the last entry in the definition of $\delta$,
which gives $|h|\le\frac{2\varepsilon}{3(|c|+1)C_1(C_0+\|D\|+1)}$.

Adding, $|\psi(w(y))-\psi(w(x_0))-D_\psi(h)-\frac12B_\psi(h,h)|\le\varepsilon|h|^{2}$ for
$|h|<\delta$. Since $\varepsilon>0$ was arbitrary, $(D_\psi,B_\psi)$ is a second-order jet
of $\psi\circ w$ at $x_0$. This is (1). (The cubic term $T_3$ is where the paper's
$\mathcal Xw\otimes\mathcal Xw$ term comes from: it is $\frac12c\,D(h)^{2}$ that survives,
the discrepancy $S^{2}-D(h)^{2}$ being of order $|h|^{3}$.)

\medskip\noindent\textbf{Step 3: proof of (2), and the splitting of the second jet.} By
\noderef{hgrad}, $(\mathrm{hgrad}\ \sigma\ x_0\ D)_i=D(\sigma^{i}(x_0))$, which is linear
in the $1$-jet; hence
$\mathrm{hgrad}\ \sigma\ x_0\ (tD)=\big(tD(\sigma^{i}(x_0))\big)_i=t\,\xi_0$, which is (2).
In particular $t\xi_0\ne0$, since $t>0$ and $\xi_0\ne0$.

Next, by \noderef{hhessS} and \noderef{Mfield},
\[
 \big(\mathrm{hhessS}\ \sigma\ x_0\ D'\ B'\big)_{i,j}
 =B'\big(\sigma^{i}(x_0),\sigma^{j}(x_0)\big)+\mathcal M(x_0,D')_{i,j},\qquad
 \mathcal M(x_0,D')_{i,j}=\tfrac12\Big(D'\big(D\sigma^{j}\sigma^{i}\big)+D'\big(D\sigma^{i}\sigma^{j}\big)\Big),
\]
so $\mathcal M(x_0,\cdot)$ is linear in the $1$-jet and $\mathcal M(x_0,tD)=t\,\mathcal M(x_0,D)$.
Moreover $(D\otimes D)(\sigma^{i}(x_0),\sigma^{j}(x_0))=D(\sigma^{i}(x_0))D(\sigma^{j}(x_0))
=(\xi_0)_i(\xi_0)_j$ by \noderef{hgrad}. Hence, writing
$\Xi:=\big((\xi_0)_i(\xi_0)_j\big)_{i,j}=\xi_0\otimes\xi_0$ and
$M:=\mathrm{hhessS}\ \sigma\ x_0\ D\ B$,
\[
 \mathrm{hhessS}\ \sigma\ x_0\ (tD)\ \big(tB+c\,D\otimes D\big)
 =\big(tB(\sigma^{i},\sigma^{j})+c(\xi_0)_i(\xi_0)_j\big)_{i,j}+t\,\mathcal M(x_0,D)
 =t\,M+c\,\Xi . \tag{5}
\]
This is the paper's
$\mathcal X\mathcal X^{s}(\psi\circ w)=\psi'(w)\mathcal X\mathcal X^{s}w+\psi''(w)\mathcal Xw\otimes\mathcal Xw$.

\medskip\noindent\textbf{Step 4: (2.5), and the proof of (3).} First, the trace against
$\Xi$ is the ellipticity function: by \noderef{Ellip},
\[
 \Tr\big(A(\xi_0)\,\Xi\big)=\sum_i\sum_j A(\xi_0)_{i,j}\,\Xi_{j,i}
 =\sum_i\sum_j A(\xi_0)_{i,j}(\xi_0)_j(\xi_0)_i=\mathcal E(A,\xi_0). \tag{6}
\]
Second, (2.5). By (2.3) (\noderef{NormalizedAH}) and $t>0$, $\xi_0\ne0$,
\[
 A(t\xi_0)=A(\xi_0),\qquad H(t\xi_0)=t\,H(\xi_0), \tag{7}
\]
whence, for any symmetric $X\in\mathbb R^{m\times m}$,
\[
 -t\Tr\big(A(t\xi_0)X\big)+H(t\xi_0)=-t\Tr\big(A(\xi_0)X\big)+t\,H(\xi_0)
 =t\Big[-\Tr\big(A(\xi_0)X\big)+H(\xi_0)\Big],
\]
which is (2.5); those are the two lines promised.

Now compute, using \noderef{Lop}, then (2) and (5), then (7), then linearity of the trace
and (6):
\begin{align*}
 \mathrm{Lop}\ \sigma\ A\ H\ x_0\ (tD)\ \big(tB+c\,D\otimes D\big)
 &=-\Tr\Big(A\big(\mathrm{hgrad}\ \sigma\ x_0\ (tD)\big)\,
    \mathrm{hhessS}\ \sigma\ x_0\ (tD)\ \big(tB+c\,D\otimes D\big)\Big)\\
 &\qquad+H\big(\mathrm{hgrad}\ \sigma\ x_0\ (tD)\big)\\
 &=-\Tr\Big(A(t\xi_0)\big(tM+c\,\Xi\big)\Big)+H(t\xi_0)\\
 &=-\Tr\Big(A(\xi_0)\big(tM+c\,\Xi\big)\Big)+t\,H(\xi_0)\\
 &=-t\,\Tr\big(A(\xi_0)M\big)-c\,\Tr\big(A(\xi_0)\Xi\big)+t\,H(\xi_0)\\
 &=t\Big[-\Tr\big(A(\xi_0)M\big)+H(\xi_0)\Big]-c\,\mathcal E(A,\xi_0)\\
 &=t\,\mathrm{Lop}\ \sigma\ A\ H\ x_0\ D\ B-c\,\mathcal E(A,\xi_0).
\end{align*}
Since $t>0$ we may factor $t$ out of the last line:
\[
 \mathrm{Lop}\ \sigma\ A\ H\ x_0\ (tD)\ \big(tB+c\,D\otimes D\big)
 =t\Big[\mathrm{Lop}\ \sigma\ A\ H\ x_0\ D\ B-\frac ct\,\mathcal E(A,\xi_0)\Big],
\]
which is (3), i.e. (3.3) with $c=\psi''(w(x_0))$ and $t=\psi'(w(x_0))$.

\medskip\noindent\emph{Remark on the hypothesis on $\psi$.} Only $\psi'(w(x_0))>0$ is
assumed, not the paper's global $\psi'>0$. That is all the proof uses: $t>0$ enters
exactly twice, in (7) --- where (2.3) needs a positive factor --- and in the division
by $t$ at the end. Stating it pointwise is necessary for the use made of this lemma later:
the relabelling $\psi_\lambda(s)=s+\lambda e^{s}$ has $\psi_\lambda'>0$ everywhere but its
companion corollary needs a pointwise lower bound on $\psi_\lambda''$, and keeping both
hypotheses pointwise keeps the two statements uniform. Nothing is lost, since a $\psi$
with $\psi'>0$ everywhere satisfies the hypothesis at every point.
\end{proof}
